% notes/03_non_realisable_resolution.tex
%
% Research notes — resolution of the non-realisable case in the
% universal spectral lower bound program.
%
% Main result: there is no separate "non-realisable regime" requiring
% a distinct minimax theory. The realisable theorem of notes/02
% (Theorem 2 corrected) applies uniformly to any learner, with the
% reachable hypothesis class $\Hcal_d$ replacing the ambient subspace
% $\Scal_\Acal$ throughout. The $\sigma_\Pi$ factor that motivated
% Conjecture 1 of notes/02 is an artefact of choosing the
% Pesenson-ridge construction as the witness predictor; using the
% population-risk minimiser of $\Hcal_d$ instead avoids it entirely.
%
% Conjecture 1 of notes/02 is therefore withdrawn as vacuous.
% Q1 from notes/01 (and refined in notes/02 §5) is FULLY RESOLVED.

\documentclass[11pt,a4paper]{article}

\usepackage[margin=2.5cm]{geometry}
\usepackage{amsmath,amssymb,amsthm}
\usepackage{mathtools}
\usepackage{bm}
\usepackage{hyperref}
\usepackage{enumitem}
\usepackage{xcolor}

\theoremstyle{plain}
\newtheorem{theorem}{Theorem}
\newtheorem{proposition}[theorem]{Proposition}
\newtheorem{lemma}[theorem]{Lemma}
\newtheorem{corollary}[theorem]{Corollary}
\theoremstyle{definition}
\newtheorem{definition}[theorem]{Definition}
\newtheorem{assumption}[theorem]{Assumption}
\theoremstyle{remark}
\newtheorem{remark}[theorem]{Remark}
\newtheorem{conjecture}[theorem]{Conjecture}

\newcommand{\bphi}{\boldsymbol{\phi}}
\newcommand{\by}{\mathbf{y}}
\newcommand{\bP}{\mathbf{P}}
\newcommand{\bB}{\mathbf{B}}
\newcommand{\bI}{\mathbf{I}}
\newcommand{\bM}{\mathbf{M}}
\newcommand{\bv}{\mathbf{v}}
\newcommand{\bone}{\mathbf{1}}
\newcommand{\E}{\mathbb{E}}
\newcommand{\R}{\mathbb{R}}
\newcommand{\Var}{\mathrm{Var}}
\newcommand{\Rspec}{R^2_{\mathrm{spec}}}
\newcommand{\Scal}{\mathcal{S}}
\newcommand{\Acal}{\mathcal{A}}
\newcommand{\Hcal}{\mathcal{H}}
\newcommand{\Ycal}{\mathcal{Y}}
\newcommand{\Rcal}{\mathcal{R}}
\newcommand{\frR}{\mathfrak{R}}

\title{\textbf{The non-realisable case is vacuous}\\
\large{Research notes — closure of Q1, withdraws Conjecture 1 of notes/02}}
\author{Rohan Fossé \and Gaël Pallares\\
\small{CESI LINEACT (EA 7527), Montpellier, France}}
\date{Working notes, May 2026 — \emph{not for distribution beyond co-authors}}

\begin{document}

\maketitle

\begin{abstract}
notes/01 and notes/02 left open the minimax rate in a "non-realisable
regime" where the orthogonal projection $\bP_\Scal\by$ falls outside the
learner's hypothesis class $\Hcal_d$ and a Pesenson-ridge witness with
$\sigma_T$-dependent slack must be used instead.
We show that this regime is \emph{vacuous}: the realisable theorem
(notes/02 Theorem~2, corrected) applies uniformly to any learner, with
the reachable hypothesis class $\Hcal_d$ replacing the ambient subspace
$\Scal_\Acal$ throughout the proof.  The $\sigma_T$ factor in notes/01
Theorem~2 and notes/02 Conjecture~1 is an artefact of choosing the
Pesenson-ridge construction as the witness predictor; replacing it with
the population-risk minimiser $f^\star_{\Hcal_d} := \arg\min_{f \in
\Hcal_d} L_V(f)$ —  which always exists, even when $\Hcal_d$ is
non-convex or non-linear — eliminates the $\sigma_T$ dependence entirely.
We formally withdraw Conjecture~1 of notes/02.  Q1 of notes/01 is
fully resolved: the universal lower bound (notes/01 Theorem~1) is
minimax-tight up to $\rho(\Pi)$ constants for any learner of any
form, with the Berry--Esseen rate $\Omega(M^2/\sqrt{N-n})$ of
notes/02 Theorem~3 as the matching lower bound.
\end{abstract}

% =========================================================================
\section{The puzzle and its resolution}
\label{sec:puzzle}

\subsection{Where the $\sigma_T$ entered}

In notes/01, the saturation theorem (Theorem~2) was proved by exhibiting
the Pesenson-ridge estimator
\begin{equation}
    \hat f_\tau \;:=\; \bB(\bB^\top \bP_T \bB + \tau \bI_d)^{-1} \bB^\top \bP_T \by
    \label{eq:pesenson-ridge-recall}
\end{equation}
on an orthonormal basis $\bB$ of $\Scal_\Acal$, with reconstruction
error bounded by $\tau^2/(\sigma_T + \tau)^2 \cdot \|\bP_\Scal\by\|^2 +
\|\bP_T \bP_{\Scal^\perp}\by\|^2/(\sigma_T + \tau)^2$
(notes/01 Lemma~3).  The choice $\tau \to 0$ recovers OLS-on-$T$ with
slack scaling as $1/\sigma_T^2$ in expectation; the choice
$\tau = \sigma_T \sqrt{d/n}$ balances bias and variance but
preserves the $1/\sigma_T$-type dependence.

notes/02 corrected this for the \emph{realisable} case (where
$\bP_{\Scal_\Acal}\by \in \Hcal_d$) by showing that the ideal projector
itself serves as a witness with slack $O(\rho M^2/\sqrt{n})$ —
\emph{no} $\sigma_T$ — and proved this rate minimax-tight via Berry--Esseen
(notes/02 Theorem~3).  Conjecture~1 of notes/02 retained the
$\sigma_T$-dependent rate $\Theta(M^2/(\sigma_\Pi\sqrt{N-n}))$ for the
"non-realisable" case ($\bP_{\Scal_\Acal}\by \notin \Hcal_d$), pointing
to Bauer--Pereverzev 2005 minimax inverse-problem theory.

\subsection{The resolution}

The non-realisable case is vacuous.  More precisely, the realisable
theorem of notes/02 applies uniformly to \emph{every} learner, with the
following two substitutions:
\begin{itemize}[label=$\bullet$,leftmargin=*]
\item Replace the orthogonal projector $\bP_{\Scal_\Acal}\by$ by the
    \emph{population-risk minimiser}
    \begin{equation}
        f^\star_{\Hcal_d} \;:=\; \arg\min_{f \in \Hcal_d} L_V(f).
        \label{eq:fstar-hd}
    \end{equation}
    This minimiser always exists by compactness (under standard truncation
    of $\Hcal_d$) and reduces to $\bP_{\Hcal_d}\by$ when $\Hcal_d$ is a
    closed linear subspace.

\item Replace the spectral $R^2$ of $\Scal_\Acal$ by the
    \emph{learner-specific spectral $R^2$}:
    \begin{equation}
        \Rspec(\Hcal_d, \by) \;:=\; 1 - \frac{L_V(f^\star_{\Hcal_d})}{\Var(\by)}.
        \label{eq:rspec-hd}
    \end{equation}
    For closed linear $\Hcal_d$, this is exactly the projection-$R^2$ of
    $\by$ on $\Hcal_d$.  For non-linear $\Hcal_d$ (gradient-boosted trees,
    MLP with weight constraints, $k$-NN), it is the fraction of
    target variance explainable in principle by some predictor in
    $\Hcal_d$.
\end{itemize}
Under these substitutions, notes/02 Theorem~2 and Theorem~3 hold verbatim,
and the universal lower bound (notes/01 Theorem~1) is minimax-tight to
$\rho(\Pi)$ constants for \emph{every} learner — linear, non-linear,
convex or non-convex hypothesis class.  No $\sigma_T$ ever appears.

% =========================================================================
\section{Formal restatement}
\label{sec:formal}

\begin{definition}[Learner-specific reachable risk]
\label{def:reachable-risk}
For any hypothesis class $\Hcal_d \subseteq \R^N$ closed under truncation
to $[-M, M]$, define
\begin{equation}
    f^\star_{\Hcal_d} \;:=\; \arg\min_{f \in \Hcal_d} L_V(f),
    \qquad
    \Rspec(\Hcal_d, \by) \;:=\; 1 - \frac{L_V(f^\star_{\Hcal_d})}{\Var(\by)}.
    \label{eq:learner-rspec}
\end{equation}
For closed linear $\Hcal_d$, $f^\star_{\Hcal_d} = \bP_{\Hcal_d}\by$ and
$\Rspec(\Hcal_d, \by) = \|\bP_{\Hcal_d}\by\|^2/\|\by\|^2$, matching
Definition~3 of notes/01.
\end{definition}

\begin{theorem}[Universal upper bound, learner-specific form — supersedes
notes/02 Theorem~2]
\label{thm:learner-specific-upper}
For any learner $\Acal$ with hypothesis class $\Hcal_d$ closed under
truncation to $[-M, M]$ and Rademacher complexity $\mathcal{R}_n(\Hcal_d)$,
under Assumption~A1 (bounded target), Assumption~A2 (ERM property),
and the existence of $f^\star_{\Hcal_d}$ in \eqref{eq:learner-rspec},
with probability $\geq 1 - \eta$ over $T \sim \Pi_{\mathrm{LSO}}$,
\begin{equation}
    L_{T^c}(\hat f_\Acal(T))
    \;\leq\; \big(1 - \Rspec(\Hcal_d, \by)\big) \cdot \Var(\by)
        \;+\; \rho(\Pi) \cdot \Bigg[
            2 \mathcal{R}_n(\Hcal_d)
            + 3 M^2 \sqrt{\frac{\log(2/\eta)}{2 n}}
          \Bigg]
        + O\!\left(\frac{M^2}{\sqrt{n}}\right).
    \label{eq:learner-upper}
\end{equation}
\end{theorem}

\begin{proof}
We replicate the proof of notes/02 Theorem~2 with $\bP_{\Scal_\Acal}\by$
replaced by $f^\star_{\Hcal_d}$.

\textit{Witness.}  $f^\star_{\Hcal_d} \in \Hcal_d$ by construction (it is
the argmin over $\Hcal_d$).  Its population loss is exactly
$(1 - \Rspec(\Hcal_d, \by))\Var(\by)$ by Definition~\ref{def:reachable-risk}.

\textit{In-sample concentration.}  Serfling without-replacement Hoeffding
inequality on the bounded sequence
$\{(f^\star_{\Hcal_d}(v) - y_v)^2\}_{v \in V}$ — which has population
mean $L_V(f^\star_{\Hcal_d}) = (1 - \Rspec(\Hcal_d, \by))\Var(\by)$
and supremum $\leq 4M^2$ — gives, with probability $\geq 1 - \eta/2$:
\begin{equation}
    L_T(f^\star_{\Hcal_d}) \;\leq\; (1 - \Rspec(\Hcal_d, \by))\Var(\by) + 4M^2 \sqrt{\frac{(1 - (n-1)/N)\log(2/\eta)}{2 n}}.
\end{equation}

\textit{ERM.}  By Assumption~A2, $L_T(\hat f_\Acal(T)) \leq L_T(f^\star_{\Hcal_d})$
since $f^\star_{\Hcal_d} \in \Hcal_d$ and ERM minimises $L_T$ over $\Hcal_d$.

\textit{Generalisation.}  Mohri Theorem~3.7 + Talagrand contraction
(square loss is $2M$-Lipschitz on $[-M, M]$) gives, with probability
$\geq 1 - \eta/2$:
\begin{equation}
    L_V(\hat f_\Acal(T)) \;\leq\; L_T(\hat f_\Acal(T)) + 2 \mathcal{R}_n(\Hcal_d) + 3 M^2 \sqrt{\frac{\log(2/\eta)}{2 n}}.
\end{equation}

\textit{Transduction.}  $L_{T^c} = L_T + \rho(L_V - L_T)$ from notes/01
Lemma~2.  Combining the above three displayed inequalities yields
\eqref{eq:learner-upper}.  $\square$
\end{proof}

\begin{theorem}[Universal lower bound, learner-specific form — supersedes
notes/02 Theorem~3]
\label{thm:learner-specific-lower}
Under the same setup, for any learner $\Acal$,
\begin{equation}
    \E_{T \sim \Pi_{\mathrm{LSO}}}[L_{T^c}(\hat f_\Acal(T))]
    \;\geq\; (1 - \Rspec(\Hcal_d, \by)) \cdot \Var(\by),
    \label{eq:learner-lower-expectation}
\end{equation}
and in high probability,
\begin{equation}
    \sup_{\by \in \Ycal_M} \mathbb{P}_T\Big[
        L_{T^c}(\hat f_\Acal(T)) - (1 - \Rspec(\Hcal_d, \by))\Var(\by) \geq c M^2 / \sqrt{N - n}
    \Big] \;\geq\; 1 - \eta
    \label{eq:learner-lower-hp}
\end{equation}
for a constant $c > 0$ depending only on $\eta$ (Berry--Esseen argument
of notes/02 Theorem~3 with the adversarial signal $\by^\star \in
\Hcal_d^\perp$, which exists whenever $\Hcal_d \subsetneq \R^N$).
\end{theorem}

\begin{proof}
Identical to notes/02 §3 with $\Scal_\Acal$ replaced by $\Hcal_d$
throughout.  The Berry--Esseen anticoncentration on $\|\by^\star\|^2_{T^c}$
applies whenever the adversarial signal lies in $\Hcal_d^\perp$, which
is possible whenever $\Hcal_d \subsetneq \R^N$ (and the latter excludes
only trivial fully-expressive learners for which the bound is vacuous
anyway).  $\square$
\end{proof}

% =========================================================================
\section{Why the $\sigma_T$ disappeared}
\label{sec:why}

The Pesenson-ridge estimator of notes/01 uses the sampling-restricted
inverse $(\bB^\top \bP_T \bB + \tau \bI_d)^{-1}$, whose operator norm is
$1/(\sigma_T + \tau)$.  This is what produces the $1/\sigma_T$
dependence in notes/01 Theorem~2.

The population-risk minimiser $f^\star_{\Hcal_d}$ does \emph{not} use
any sampling-restricted inverse: it is defined directly on the
population $V$, with $L_V$ as the objective.  The slack of
\eqref{eq:learner-upper} is therefore controlled by classical empirical
process theory (Rademacher + Hoeffding--Serfling), which has nothing
to do with the sampling quality $\sigma_T$ of any specific subspace.

The Pesenson-ridge construction is a \emph{specific algorithm}, not a
fundamental rate.  Using a different ERM-flavoured algorithm (LightGBM,
kernel ridge, neural network) avoids the $\sigma_T$ cost entirely.

\begin{remark}[Computational accessibility of $f^\star_{\Hcal_d}$]
\label{rem:fstar-accessible}
For closed linear $\Hcal_d$, $f^\star_{\Hcal_d} = \bP_{\Hcal_d}\by$ requires
knowing $\by$ on all of $V$ — which the learner does not have.  This
matters for the \emph{algorithm}, not for the \emph{minimax rate}: the
rate measures what is achievable in principle, and the ERM learner
\emph{approximately recovers $f^\star_{\Hcal_d}$} from $\by|_T$ via
empirical process theory (Mohri 2018), incurring only the Rademacher
slack already accounted for in \eqref{eq:learner-upper}.  No knowledge
of $\by$ on $T^c$ is required; the witness $f^\star_{\Hcal_d}$ enters
the proof, not the algorithm.
\end{remark}

\begin{remark}[Connection to the MLST proof]
\label{rem:mlst-proof}
The MLST paper (\texttt{materials-applicability-bound/applicability\_bound.tex},
Section 5.2, Part (c)) implicitly uses
$f^\star_{\Hcal_d} = \bP_{\Scal_{\mathrm{comp}}}\by$ as its witness, with
$\Hcal_d \supseteq \Scal_{\mathrm{comp}}$ by Assumption (the hypothesis
class subsumes the linear projector on the compositional subspace).
There is no Pesenson-ridge ambiguity in that paper, and no $\sigma_T$
factor: the proof structure of Theorem~\ref{thm:learner-specific-upper}
above is identical to MLST Part (c).
\end{remark}

% =========================================================================
\section{Withdrawal of notes/02 Conjecture~1}
\label{sec:withdrawal}

We formally withdraw Conjecture~1 of \texttt{notes/02\_minimax\_saturation.tex}
(\S 3, Non-realisable case).  Its stated form
\begin{equation*}
    \mathfrak{R}_n^{\mathrm{non-real}}(\Ycal_{\Scal,M}, \eta)
    \;=\; \Theta\!\left(\frac{M^2}{\sigma_\Pi(\Hcal_d) \cdot \sqrt{N-n}}
        + \frac{M^2}{\sqrt{N-n}}\right)
\end{equation*}
contains a spurious $1/\sigma_\Pi$ factor on the first term, traced to
the conflation of (i) the Pesenson-ridge construction's
operator-norm slack and (ii) the genuine minimax rate of the
non-realisable problem.

Under the resolution of §\ref{sec:puzzle}, the correct rate is
\begin{equation}
    \mathfrak{R}_n(\Hcal_d, M, \eta)
    \;=\; \Theta\!\left(\frac{M^2}{\sqrt{N-n}} + \rho(\Pi) \mathcal{R}_n(\Hcal_d)\right),
    \label{eq:correct-rate-hd}
\end{equation}
matching the realisable rate of notes/02 Theorem~2 with $\Hcal_d$
replacing $\Scal_\Acal$.  The $\sigma_\Pi$ term is absent.  The
"forthcoming notes/03" referenced in Bauer--Pereverzev paragraph of
notes/02 §3 is the present document; the Bauer--Pereverzev minimax
inverse-problem theory turns out not to be needed for the program,
because our setting is not the noisy linear inverse problem they treat
but the deterministic empirical-risk problem with random
test-set splitting, for which classical statistical learning theory
(Rademacher + Hoeffding--Serfling) is exactly the right tool.

% =========================================================================
\section{Status of the program after notes/01–03}
\label{sec:status}

The universal spectral lower bound program is, at the formal-proof level,
complete.

\begin{itemize}[label=$\bullet$, leftmargin=*]
\item \textbf{Lower bound (notes/01 Theorem~1).}  Universal,
    learner-independent, exact in expectation: $\E[L_{T^c}(\hat f)] \geq
    (1 - \Rspec(\Hcal_d, \by)) \Var(\by)$, with $\Hcal_d$ the learner's
    hypothesis class (Definition~\ref{def:reachable-risk}).  \emph{Proved.}

\item \textbf{Upper bound (Theorem~\ref{thm:learner-specific-upper}
    above; supersedes notes/02 Theorem~2).}  ERM learner achieves the
    lower bound with slack $\rho(\Pi) \cdot [2\mathcal{R}_n(\Hcal_d) +
    O(M^2/\sqrt{n})]$.  \emph{Proved.}

\item \textbf{Minimax tightness
    (Theorem~\ref{thm:learner-specific-lower}; supersedes notes/02
    Theorem~3).}  The slack is minimax-tight up to $\rho(\Pi)$
    constants; Berry--Esseen forces $\Omega(M^2/\sqrt{N-n})$ even for
    oracle estimators.  \emph{Proved.}

\item \textbf{Non-realisable regime.}  Vacuous.  \emph{Closed.}

\item \textbf{Corollaries C1, C2, C3.}  C1 = MLST Theorem~1 (corollary of
    Theorem~\ref{thm:learner-specific-upper} with $\Hcal_d = \Scal_{\mathrm{cat}}$
    vs $\Scal_{\mathrm{comp}}$).  C2 (negative transfer) and C3 (active
    learning) inherit the same framework; full corollaries to be written
    inside the submission-ready manuscript.
\end{itemize}

The remaining open questions from notes/01 §6 and notes/02 §5 reduce to
two non-foundational ones:
\begin{enumerate}
\item \textbf{Sharp $\rho(\Pi)$ constant} (notes/02 §5, item 5).  A
    Fano-style multi-point argument should give tight constants.
    Mathematical exercise, not foundational.
\item \textbf{Importance-weighted variant for non-uniform protocols}
    (notes/01 §6, item 3).  Required for active-learning corollary C3.
    Mechanical, not foundational.
\end{enumerate}

\begin{remark}[Implication for the manuscript]
\label{rem:manuscript-implication}
The submission-ready manuscript for
\texttt{structural-bounds-framework} can be drafted entirely from
notes/01 (Theorem 1 + Step-1/2/3 proof structure) and the corrected
upper-bound statement of Theorem~\ref{thm:learner-specific-upper} above.
Neither the Pesenson-ridge estimator nor any $\sigma_T$ factor needs
to appear in the manuscript.  This dramatically simplifies the writeup:
the bound is a clean realisable-case empirical-process result with a
spectral interpretation, not an inverse-problem result.
\end{remark}

% =========================================================================
\section*{Acknowledgements (working notes)}

This resolution emerged on 2026-05-24 from a careful per-component
analysis of the $\sigma_T$-dependence in the Pesenson-ridge slack
of notes/01 Theorem~2.  The key observation — that the
population-risk minimiser $f^\star_{\Hcal_d}$ serves as a witness
predictor with no $\sigma_T$ cost — was triggered by re-reading
the proof of MLST Theorem~1(c) and noticing that it never uses
the spectral nature of the witness $\bP_{\Scal_{\mathrm{comp}}}\by$.

Conjecture~1 of notes/02 is withdrawn.  The Bauer--Pereverzev 2005
pointer in that note's §3 is, in retrospect, a wrong-domain reference:
their theory applies to noisy inverse problems on continuous Hilbert
spaces, while our setting is the empirical-risk problem with random
sub-sampling on a finite graph — a different mathematical object.

\end{document}
